\section{The quadratic variation and its isogeny exceptions}
\label{sec:geometry}

Quaternionic multiplication will be recorded by two independent Hodge
vectors in a rank-five quadratic local system.  We first construct this
system on a fixed fine-level curve and determine its monodromy.  For two
curve arguments, failure of independent monodromy forces a polarized
isogeny relation.  We show that only finitely many such relations can
dominate the fixed curves.  This finiteness will allow the auxiliary level
in \Cref{sec:height} to be chosen before any target point is considered.

\subsection{Fixed curves and finite covers}

\begin{lemma}[Passage to fixed covers]\label{lem:finite-covers}
Let $C\subset\Acal_2$ be a reduced irreducible closed curve over $\Qbar$.
There is a finite surjective map $\pi:S\to C$ from a smooth geometrically
connected curve over a number field, together with a principally polarized abelian scheme
$f:\mathcal A\to S$ whose coarse moduli map is $\pi$.  One may obtain $S$
by normalization and a fixed full symplectic level $N\geq3$.

If infinitely many points of $C$ have the endomorphism algebra in
\Cref{thm:main}, then infinitely many points of $S$ have that property.
Deleting finitely many points of $S$, or making a further fixed finite
surjective cover and selecting a dominating connected component, preserves
this assertion.  Once all fixed data are defined over $K$, degrees and ample
logarithmic heights on the smooth completions satisfy
\begin{equation}\label{eq:fixed-cover-bounds}
 [K(\pi(t)):K]\leq[K(t):K]\leq D[K(\pi(t)):K],\qquad
 c^{-1}h_C(\pi(t))-c\leq h_S(t)\leq c h_C(\pi(t))+c,
\end{equation}
with fixed constants $D,c$.  The same statements hold for any further fixed
cover.  In particular, a degree-dependent polynomial height bound obtained
on one selected lift of every target descends to the original targets.
\end{lemma}

\begin{proof}
The full level moduli space is a fine moduli scheme with a universal abelian
scheme \cite[Chapter~I, Theorem~1.4]{Chai1985}; for the complex moduli
interpretation see also \cite[Theorem~4.6 and Corollary~4.7]{MilneModuli2013}.
The finite change-of-level group acts on it, and its quotient is the
coarse moduli space, so the forgetful map is finite.  We fix the
symplectic pairing component after adjoining the relevant roots of unity.
Normalize a component of its
pullback dominating $C$.  Normal curves in characteristic zero are smooth;
the resulting map is finite and surjective.  All the schemes and morphisms
in this construction descend to a number field.  Fibres of a finite map are
finite, so an infinite set downstairs has infinitely many lifts.  For a
cover with several components, at least one component contains infinitely
many lifts; a dominating component of a finite cover of an irreducible
curve is itself surjective.

After normalizing the base, the finite maps extend to finite maps between
smooth projective completions.  For such a map of generic degree $D_0$,
the degree formula for a fibre gives
$[K(t):K(\pi(t))]\leq D_0$.  Over the finitely many singular points of
the original coarse curve, the fixed finite map has a fixed finite set
of preimages; enlarge $D$ to bound their residue-field extensions as
well.  Thus the constant $D$ in \eqref{eq:fixed-cover-bounds} is uniform,
but need not equal the generic degree of the map to the singular curve.

Functoriality of heights identifies a height of a pulled-back ample
divisor with the height downstairs up to a bounded function.  Any two
ample divisors on a fixed projective curve bound each other's heights
linearly: sufficiently large positive integral multiples of either minus
the other are ample and have heights bounded below.  This proves
\eqref{eq:fixed-cover-bounds}.  Replacing the fixed field $K$ changes
degrees by at most a fixed factor.  None of these constants depends on the
point.
\end{proof}

Assume the target locus is infinite, as will be appropriate for the final
contradiction, and choose $S$ as in \Cref{lem:finite-covers}.  We write
$\bar S$ for its smooth projective completion.  After any required finite
deletions, choose a target $t_0\in S(\Qbar)$ and enlarge the fixed number
field $K$ so that $t_0\in S(K)$.  A rational function $x\in K(S)$ can be
chosen with a simple zero at $t_0$.  On a neighbourhood of $t_0$, it is an
\'etale parameter; its inverse germ through $t_0$ is always the germ used
below.  Other zeros of $x$ do not specify this germ.  Rational frames and
affine coordinates may be made regular on this neighbourhood by further
finite deletions away from $t_0$.  The finitely many conjugates of these
fixed data will also be treated as fixed data.  No property of the
reduction of $t_0$ is imposed here.

\subsection{A rank-five operator model}

Let $\mathbb H_{\Z}=R^1f_*\Z$ on $S(\C)$, and let $\psi$ be its alternating
polarization.  The corresponding rational local system is
$\mathbb H=\mathbb H_{\Z}\otimes\Q$.  The adjoint of an operator for $\psi$
is denoted by $T^\dagger$.  We use the weight-zero local system
\begin{equation}\label{eq:quadratic-system}
 \mathbb V=\{T\in\End(\mathbb H):T^\dagger=T,\ \Tr(T)=0\},
 \qquad \pair{T}{U}=\tfrac14\Tr(TU),\qquad q(T)=\pair{T}{T}.
\end{equation}
Its natural integral lattice is the intersection with
$\End(\mathbb H_{\Z})$.  The rational form in
\eqref{eq:quadratic-system} is integral and unimodular on this lattice
after inverting $2$.  Using another fixed commensurable integral lattice
only enlarges the finite set of excluded auxiliary primes.

\begin{lemma}[The quadratic model and the quaternionic plane]
\label{lem:qm-plane}
The local system $\mathbb V$ has rank five, a nondegenerate form of real
signature $(3,2)$, and Hodge numbers $(1,3,1)$ in types
$(1,-1),(0,0),(-1,1)$.  It is the primitive part of $R^2f_*\Q(1)$.
Conjugation of operators gives a surjective morphism of algebraic groups
\begin{equation}\label{eq:gsp-to-so}
 \rho:\GSp(\mathbb H,\psi)\longrightarrow\SO(\mathbb V,q),
 \qquad g:T\longmapsto gTg^{-1},
 \qquad \ker\rho=\mathbb G_m.
\end{equation}
Its restriction to $\Sp_4$ is the spin covering, with kernel
$\{\pm\id\}$.

At a point $t$ whose full geometric endomorphism algebra is an indefinite
quaternion division algebra, the trace-free Rosati-symmetric
endomorphisms form a rational two-plane
$E_t\subset\mathbb V_t$ of Hodge vectors.  The restriction of $q$ to $E_t$
is positive definite.  This plane is spanned by two actual integral
trace-free symmetric endomorphisms.  If $\Lambda_t=\mathbb V_{\Z,t}$, then
$L_t=E_t\cap\Lambda_t$ is saturated in $\Lambda_t$; a selected generating pair need not
generate this saturated lattice.
\end{lemma}

\begin{proof}
In a symplectic basis, the matrix of $\psi$ is
$J=\left(\begin{smallmatrix}0&I_2\\-I_2&0\end{smallmatrix}\right)$.
The equations $T^tJ=JT$ and $\Tr(T)=0$ give
\[
 T=\begin{pmatrix}A&\beta J_2\\ \gamma J_2&A^t\end{pmatrix},
 \qquad A=\begin{pmatrix}a&b\\c&-a\end{pmatrix},
 \qquad J_2=\begin{pmatrix}0&1\\-1&0\end{pmatrix}.
\]
One checks that $T^2=(a^2+bc-\beta\gamma)I_4$.  Thus
\[
 q(T)=a^2+bc-\beta\gamma,
\]
which has signature $(3,2)$ and is nondegenerate over $\Z[1/2]$.
The alternating form $(u,v)\mapsto\psi(Tu,v)$ identifies the
self-adjoint operators with alternating two-forms.  The identity operator
corresponds to $\psi$, and its orthogonal complement for the trace form
is the primitive part.  Taking account of the polarization twist gives
$\mathbb V\simeq(\bigwedge^2\mathbb H)(1)_{\mathrm{prim}}$ and the stated
Hodge numbers.

Conjugation by a similitude preserves self-adjointness and the trace
pairing.  Its determinant on $\mathbb V$ is $1$, since $\GSp_4$ is
connected.  The induced Lie algebra map
$\mathfrak{sp}_4\to\mathfrak{so}_5$ is nonzero.  Simplicity and equality
of dimensions make it an isomorphism.  Consequently the induced map on
the adjoint groups is an isomorphism; these are the split adjoint groups
of types $C_2$ and $B_2$.  Its kernel on $\GSp_4$ is exactly the scalar
torus, proving \eqref{eq:gsp-to-so} and the assertion about the spin
covering.  These statements are over $\Q$, not merely over $\C$.

Put $D=\End^0(\mathcal A_t)$.  The polarization induces the positive
Rosati involution on $D$.  Over $\R$, the identification
$D\otimes\R\simeq M_2(\R)$ turns a positive involution into the adjoint
for a positive definite form, and hence into transpose in a suitable
basis.  Its fixed subspace has dimension three.  The cohomological
trace is nonzero on the identity, so its kernel in that subspace has
dimension two.  For a nonzero symmetric endomorphism $u$, Rosati
positivity gives $\Tr(uu^\dagger)=\Tr(u^2)>0$.  Thus $q$ is positive
definite on this trace-free subspace.

The rational realization of endomorphisms on $H^1$ is faithful and
identifies endomorphisms with rational Hodge endomorphisms
\cite[Theorem~4.1]{MilneShimura1995}.  In particular
these vectors are Hodge vectors in $\mathbb V_t$.  Clearing denominators
in a rational basis gives integral endomorphisms with the same symmetry
and trace conditions.  Equivalently one can first apply
$u\mapsto u+u^\dagger$ and then
$v\mapsto4v-\Tr(v)\id$ to integral endomorphisms; the principal
polarization makes the first operation integral.  Finally, intersection
of an integral lattice with a rational subspace is saturated.
\end{proof}

Write $\Vcal_{\mathrm{dR}}$ for the algebraic de Rham realization, with its
Gauss--Manin connection and trace form.  The connection is the one
constructed from the relative de Rham complex
\cite[Theorem~1]{KatzOda1968}; its functoriality and compatibility with
cup products follow from that construction.  Algebraic--analytic de Rham
comparison \cite[Theorem~1$'$]{Grothendieck1966DeRham}, applied
to the relative complexes, identifies its horizontal local system with
the Betti realization; this relative interpretation is recorded in
\cite[Section~0.0]{Katz1970}.  Actual endomorphisms have
compatible Betti, de Rham, and \'etale realizations, and their trace
pairings are rational.  For a polarized isogeny
$\varphi:A\to A'$ with $\varphi^*\lambda'=m\lambda$, cohomological
pullback $\varphi^*:H^1(A')\to H^1(A)$ has multiplier $m$ and reverses
direction.  Our covariant transport on $H^1$ is
\[
 T_\varphi=m(\varphi^*)^{-1}:H^1(A)\longrightarrow H^1(A'),
\]
which again has multiplier $m$.  We put
$I_\varphi(u)=T_\varphi uT_\varphi^{-1}$ in the operator model.  This is
the isometry in \eqref{eq:gsp-to-so} in the direction from $A$ to $A'$;
on actual endomorphisms it realizes
$\alpha\mapsto\varphi\alpha\varphi^{-1}$.  These conventions also apply
to polarized quasi-isogenies with rational multiplier and commute with
composition.  Thus an inverse isogeny of multiplier $m$ composed with
Frobenius of multiplier $p$ has multiplier $p/m$ on this covariant
weight-one realization.  For an isogeny, $\deg\varphi=m^2$.

Under a flat trivialization, the period is the line of type $(1,-1)$ in
the smooth projective quadric
\begin{equation}\label{eq:period-quadric}
 Q_V=\{[z]\in\mathbb P(V_{\C}):q(z)=0\},\qquad \dim_{\C}Q_V=3.
\end{equation}
A rational vector is of type $(0,0)$ if and only if it is orthogonal to
this line: its reality supplies orthogonality to the conjugate line as
well.  In particular, the period at a quaternionic target is orthogonal
to $E_t$.

\subsection{Monodromy and rational isogenies}

\begin{proposition}[Single and product monodromy]\label{prop:monodromy}
For the family over any fixed smooth finite cover of a Hodge-generic
curve in $\Acal_2$, the connected algebraic monodromy on $\mathbb H$ is
$\Sp_4$, and that on $\mathbb V$ is $\SO(V,q)$.  The same holds after
pullback by a dominant morphism from an irreducible complex algebraic
variety, on a smooth open of a resolution.

Let $S_1,S_2$ be two such curves, and let an irreducible complex algebraic
variety $W$ map dominantly to both.  On the sum of the two pulled-back
weight-one systems, either the connected monodromy is
$\Sp_4\times\Sp_4$, or the image of $W$ in the product of the two moduli
curves is contained in a dominating component of a polarized isogeny
correspondence.  In the latter case the corresponding local period
relation is given by a fixed element of $\GSp_4(\Q)$ with positive
multiplier.
\end{proposition}

\begin{proof}
The generic Mumford--Tate group of a Hodge-generic curve is $\GSp_4$.
The connected monodromy normality theorem for a polarizable geometric
variation on a smooth algebraic base places connected monodromy normally
in its derived group \cite[\S5, Theorem~1]{Andre1992}.  Since $\Sp_4$ is
$\Q$-simple, the connected monodromy is either trivial or $\Sp_4$.
In the first case a finite cover kills monodromy, and the theorem of the
fixed part makes the variation constant; the moduli map would be
constant.  This is impossible.  This consequence is also recorded in
\cite[\S6, Remark~1]{Andre1992}.  A dominant pullback has the same generic
Mumford--Tate group and still has a nonconstant moduli map, so the argument
also applies there.  The assertion for $\mathbb V$ follows from
\eqref{eq:gsp-to-so}.

For the two systems let $M$ be their generic Mumford--Tate group and
$H$ their connected monodromy, after fixing rational symplectic bases
on a simply connected open.  Both projections of $H$ are $\Sp_4$.
Put $G=\operatorname{PGSp}_4$.  The algebraic Goursat lemma and simplicity
give two possibilities for the image $\bar H\subset G\times G$:
it is the product, or it is the graph of a $\Q$-defined automorphism
$\alpha$ of $G$.  Indeed the kernels of the two projections are normal;
in an adjoint simple group a proper such kernel is trivial.  In the
product case dimensions give $H=\Sp_4\times\Sp_4$.

In the graph case, every automorphism of the split adjoint group of type
$C_2$ is inner, since its Dynkin diagram has no automorphism.  As $G$ is
adjoint, a $\Q$-defined inner automorphism is conjugation by an element
$\bar g\in G(\Q)$.  The central exact sequence
\[
 1\longrightarrow\mathbb G_m\longrightarrow\GSp_4
 \longrightarrow G\longrightarrow1
\]
and $H^1(\Q,\mathbb G_m)=0$ show that $\bar g$ lifts to
$g\in\GSp_4(\Q)$.  Conjugate the second factor by this lift, so that
$\bar H$ becomes the diagonal.  Normality of $H$ in $M$ forces the
adjoint image of $M$ into the normalizer of this diagonal.  That
normalizer is the diagonal itself: a normalizing pair $(g_1,g_2)$ has
$g_2^{-1}g_1$ central in the adjoint group.

The two Hodge homomorphisms therefore agree after passage to $G$.
Their discrepancy on the weight-one spaces is a real algebraic
character of the Deligne torus
$\operatorname{Res}_{\C/\R}\mathbb G_m$ into the scalar $\mathbb G_m$.
Such a character is a power of the norm.  The two weights are equal,
so restriction to real scalars makes this power zero.  The rational
map $g$ thus identifies the Hodge structures themselves.  Its
similitude multiplier is positive because it compares the positive
definite forms supplied by their polarizations and complex structures.

This argument holds at very general points of the chosen open;
holomorphicity extends the resulting fixed period relation across it.
Write $\mu>0$ for the multiplier of
$g:H^1(A_1)\to H^1(A_2)$.  The homology map in the direction from $A_1$
to $A_2$ is
\[
 F=\mu(g^{-1})^\vee:H_1(A_1,\Q)\longrightarrow H_1(A_2,\Q).
\]
It has multiplier $\mu$ and commutes with the complex structures.
Choose an integer $n>0$ such that $nF$ is integral.  It induces a
holomorphic isogeny $\varphi:A_1\to A_2$, which is algebraic since the
tori are projective.  Its multiplier is the positive integer
$m=n^2\mu$: integrality follows by evaluating its polarization identity
on a pair of integral homology vectors of pairing $1$.
Moreover $\varphi^*=n\mu g^{-1}$, so the covariant cohomological
transport defined above is
\[
 T_\varphi=m(\varphi^*)^{-1}=ng.
\]
Thus $\varphi$ realizes the original projective rational representative.

A fixed rational similitude defines a
closed algebraic Hecke correspondence, finite over either moduli
factor.  Concretely, for a fixed integral multiplier $m$, its kernels
are among the finite subgroup schemes of $A[m]$, and the principal
polarizations and fixed level structures give the corresponding finite
moduli problem.  Closure of the local relation consequently contains
the image of $W$ in one such correspondence.  Its projections to both
curves are dominant by hypothesis.  Finite central twists of the local
systems can be killed by a finite cover and do not change this
conclusion about their moduli points.
\end{proof}

\subsection{Only finitely many dominating correspondences}

Here the finiteness concerns whole correspondences dominating two fixed
curves.  Individual isolated isogenous pairs need not belong to the
finite list.

\begin{proposition}[Finitely many isogeny exceptions]
\label{prop:finite-correspondences}
Fix finitely many closed irreducible Hodge-generic curves in
$\Acal_2(\C)$.  For every ordered pair of these curves there are finitely
many rational positive projective symplectic similitudes carrying a
lift of the first onto a lift of the second, modulo integral changes
of lifts.  Consequently all dominating polarized isogeny
correspondences between this list of curves are contained in a finite
union of closed curves in their pairwise products.  There is a fixed
finite set of positive integers $m$ such that every pair on this union
is joined by a polarized isogeny of one of these multipliers, of degree
$m^2$.  Refining the fixed level does not enlarge the required set of
underlying multipliers.
\end{proposition}

\begin{proof}
We first show that the real stabilizer of a lifted curve is discrete.
Finite area of the curve quotient will then make its integral stabilizer
a subgroup of finite index.  Maps between two lifts consequently have
only finitely many cosets, and rational representatives will give the
fixed isogeny multipliers.

Let $\mathfrak H_2$ be Siegel space, let
$G_{\R}=\operatorname{PGSp}_4(\R)^+$ be its effective connected group of
holomorphic isometries, and fix a torsion-free principal congruence
subgroup $\Gamma$ at fine level.  For any curve in the list, choose an
irreducible analytic component $X$ of its full inverse image in
$\mathfrak H_2$.  It is a closed analytic curve.  Its normalization
may be used wherever a smooth curve is needed.  The integral stabilizer
\[
 \Gamma_X=\{\gamma\in\Gamma:\gamma X=X\}
\]
contains the image of the monodromy of the normalization of the
corresponding fine-level curve.  In fact this image is the deck group
on the lifted normalization over the smooth part.  By
\Cref{prop:monodromy}, it is Zariski dense in the real adjoint group.

We verify that the real stabilizer
\[
 H_X=\{g\in G_{\R}:gX=X\}
\]
is discrete.  It is closed: if $g_n\to g$ with $g_nX=X$, closedness of
$X$ gives $gX\subset X$, and applying the same argument to $g_n^{-1}$
gives equality.  Thus $H_X$ is a Lie subgroup.  Its Lie algebra is
invariant under $\Gamma_X$.  Stabilizing a fixed real vector subspace
of the adjoint representation is an algebraic condition, so Zariski
density makes this Lie algebra an ideal of
$\mathfrak{sp}_4(\R)$.  Simplicity leaves zero or the full Lie algebra.
The latter would give $H_X=G_{\R}$, forcing $X=\mathfrak H_2$ by
transitivity.  The dimensions differ, so $H_X$ is discrete.

Discreteness alone does not bound the number of cosets of $\Gamma_X$.
For this we use that $\Gamma_X\backslash X$, with normalization understood,
has finite area for the metric induced from an invariant Hermitian
metric on $\mathfrak H_2$.  To check both the quotient and the effect
of singularities, let $S_X$ be the normalization of the corresponding
closed fine-level curve and form
\[
 Y=S_X\mathbin{\times}_{\Gamma\backslash\mathfrak H_2}\mathfrak H_2.
\]
The normalization map from $S_X$ to its closed fine-level image is
finite, so base change makes $Y\to\mathfrak H_2$ finite, with image
the full inverse image of that curve.  On the other hand,
$Y\to S_X$ is a covering and locally a biholomorphism, so $Y$ is
smooth.  The finite map to that inverse image is an isomorphism above
the smooth locus of the curve; hence it is its normalization.

The component $Y_0$ normalizing $X$ has stabilizer
$\Gamma_X$; this group acts transitively on every fibre over $S_X$, since a deck
transformation carrying two points of the same fibre to one another
preserves their common connected component, and
$\Gamma_X\backslash Y_0\simeq S_X$.  It therefore suffices to bound the
pullback metric at the finitely many punctures of $S_X$.

For completeness, any holomorphic map from the unit disk to
$\mathfrak H_2$ contracts its invariant metric up to a fixed constant
relative to the Poincar\'e metric.  Realize $\mathfrak H_2$ as a bounded
domain.  An automorphism moving the image of the origin to a fixed
origin, followed by the Cauchy derivative estimate in bounded
coordinates, bounds the derivative at zero uniformly.  Disk
automorphisms give the pointwise assertion.  Apply it on the universal
cover of a punctured parameter disk.  The descended upper bound is a
constant multiple of
\[
 \frac{|du|^2}{|u|^2\log^2(1/|u|)},
\]
whose area on $0<|u|<\varepsilon$ is finite.  On the remaining compact
part the period map is holomorphic and has finite area.  Zeros of its
derivative or singular branches only decrease this upper bound.

We now deduce the precise group-theoretic consequence
\begin{equation}\label{eq:stabilizer-index}
 [H_X:\Gamma_X]<\infty.
\end{equation}
The action of $G_{\R}$ on $\mathfrak H_2=G_{\R}/K_{\mathrm{max}}$ is
proper, since $K_{\mathrm{max}}$ is compact.  Therefore the discrete
group $H_X$ acts properly on $X$, with finite point stabilizers.  Choose
a regular point $z\in X$ and a sufficiently small $H_{X,z}$-invariant
disk $D\subset X$ of positive induced area such that $gD\cap D$ is
empty for $g\notin H_{X,z}$.  Distinct double cosets in
$\Gamma_X\backslash H_X/H_{X,z}$ give disjoint images of translates of
$D$ in $\Gamma_X\backslash X$.  Each translate embeds there: a further
identification would give a nonidentity element of the torsion-free
group $\Gamma_X$ conjugate to an element of the finite group $H_{X,z}$.
All these disks have the same positive area.  Finite total area thus
bounds the number of double cosets, and finiteness of $H_{X,z}$ bounds
the number of left cosets.  This proves \eqref{eq:stabilizer-index}.

Choose lifts $X_1,X_2$ for two curves.  If a rational positive projective
similitude $g_0$ carries $X_1$ onto $X_2$, the set of all real such
maps is $g_0H_{X_1}$.  By \eqref{eq:stabilizer-index} this set has only
finitely many classes modulo right multiplication by
$\Gamma_{X_1}$.  Every class containing a rational map admits a rational
representative.  Integral changes of the chosen lifts add no classes:
at coarse level all lifts are conjugate under the full integral
symplectic group; at the chosen fine level there are only finitely
many resulting classes.  Thus there are finitely many rational maps
for the whole fixed list.

A dominating component of a Hecke correspondence determines a local
equality between a rational translate of one lifted curve and another
lifted curve.  Both are irreducible closed analytic curves, so this
local equality extends to equality of those lifts.  It is therefore
represented by one of the finite maps just obtained.  Lift each
projective rational representative to $\GSp_4(\Q)$, using
$H^1(\Q,\mathbb G_m)=0$, and clear denominators on the integral homology
lattices.  Its multiplier becomes a positive integer $m$.  This gives an isogeny of multiplier $m$ on the local correspondence.
To retain that multiplier after taking closure, work over the fixed
fine moduli space in characteristic zero.  Subgroups
$K\subset A[m]$ that are maximal isotropic for the Weil pairing and
have order $m^2$ form a finite \'etale scheme over this space:
\'etale-locally the torsion is a constant finite group.  Each such
subgroup gives the quotient $B=A/K$ with its principal polarization
characterized by $q^*\lambda_B=m\lambda_A$.  The resulting map to
$\Acal_2\times\Acal_2$ is finite.  Indeed it is finite over the first
coarse factor, hence proper into the product because the second
factor is separated, and its fibres are finite.  Its image is
therefore closed, and every geometric point of that image admits an
isogeny of multiplier $m$.  The closure of the local correspondence
lies in this image, including at singular coarse points and points
with extra automorphisms.

Taking all representatives and the finitely many curve components of
their intersections with the pairwise products that dominate both factors
gives the required finite union.  Additional level structures change only
the finite covers of these moduli problems, not the underlying
isogenies or their multipliers.
\end{proof}

We apply \Cref{prop:finite-correspondences} to the full finite list of
absolute Galois conjugates of the original closed coarse curve $C$.
The canonical Siegel moduli space has reflex field $\Q$, and Galois
conjugation permutes its special subvarieties
\cite[Corollary~5.3, arXiv version~2]{Orr2021}.  If a conjugate of $C$ lay in a proper
special subvariety, inverse conjugation would contradict the Hodge
genericity of $C$.  Thus every curve in this list is Hodge generic.
Enlarging $K$ for centers or markings adds no coarse image to this
list.  The resulting finite set of multipliers, the chosen center,
and the rank-five lattices are fixed before the prime $\ell$ and the
marking cover of \Cref{prop:marking} are selected.  Refining level changes
the markings over these same coarse images; the exceptional isogenies
are required only between the underlying polarized abelian varieties.
Later estimates may therefore be proved on that fixed cover and
descended by \Cref{lem:finite-covers}.
